причина значної частини наших страждань --- наше незнання щодо духовної природи Всього: усіх предметів, наших власних тіл, наших стосунків з іншими людьми й нашої манери переживати час.

ми думаємо, що є розділення там, де його немає. ми відділяємо духовне життя від «нормального». ми шукаємо медитативні практики, щоб лікувати симптоми на кшталт тривоги й депресії, тоді як насправді проблема значно глибша. «хочу медитувати, щоб повернутися до свого нормального життя з меншим стресом». звісно, будь-яка причина практикувати медитацію слушна. але навіть сама причина --- це те, що слід перевершити, бо паліативна медитація не є достатньою, щоб звільнити нас від циклів страждання.

під час медитації ми просто відтворюємо ті самі розумові звички, що й решту дня. ми даємо себе захопити уму, який надмірно цінує цілі й наполегливо безперервно судить нас самих та інших.

наївно думати, що ті хвилинки медитації самі по собі вилікують хоч що-небудь. «ах, моє духовне життя --- це коли я йду до медитаційного центру»; «моє духовне життя --- це коли я практикую йогу»; «моє духовне життя --- це коли я йду до церкви». усі --- приклади прояву тієї самої помилки.

у міру того як ми дозволяємо розчинятися бар'єру, який створили між духовним життям і звичайним, медитативний стан стає дедалі присутнішим. без безперервних розумових звичок і зумовленостей ми зрештою доходимо до точки, де можемо просто дозволити собі сісти на землю без цілей і без суджень. практика стає насолодою. життя стає насолодою. це і є просвітлення.

